\documentclass[14pt,a4paper,oneside]{article}

\usepackage[utf8]{inputenc}   


\usepackage{amsmath,amsthm,amssymb,appendix,bm,mathrsfs,ctex,geometry,tcolorbox,fancyhdr,tikz,esint,appendix}


%% GB/T 7714 宏包 ---- 引文格式
\usepackage[sort&compress]{gbt7714}
\bibliographystyle{gbt7714-numerical}



% \numberwithin{equation}{section}             %%%%%% 公式编号




\geometry{left=1.7cm, right=1.7cm, top=2.54cm, bottom=2.54cm}    %%%%% 页面尺寸
\linespread{1.5}                             %%%%% 行距1.5倍
\pagestyle{fancy}



%%    设置页眉页脚
\fancyhf{}
\renewcommand{\headrulewidth}{0.15mm}
\fancyhead[C]{文献阅读报告}                   %%%%%% 页眉
\fancyfoot[C]{\thepage}                      %%%%%% 页脚 页码
\setlength{\headheight}{15pt}


%%   设置各种链接颜色为黑色
\usepackage[colorlinks,
            linkcolor=black,
            anchorcolor=black,
            citecolor=black,
            urlcolor=blue
            ]{hyperref}



%%   例子编号
\newcounter{example}
\counterwithin{example}{subsection}
\newcommand{\examplename}{Example}
\newcommand{\exampleTitle}[1]{%
\textbf{\refstepcounter{example}\examplename\theexample:#1}\par}


%%  插入图片配置
\usepackage{graphicx,subfigure}
\usepackage[]{caption2} %新增调用的宏包
% \numberwithin{figure}{section}               %%%%%% 图片编号按节

\graphicspath{ {./Figure/} }            %%%%% 设置图片文件所在目录
\DeclareGraphicsExtensions{.pdf,.svg,.jpeg,.png,.jpg}     %%%%% 支持图片格式


%%    支持tikz作图
\usepackage{tikz}
\usetikzlibrary{backgrounds,arrows}


\usepackage{pgfplots}
\pgfplotsset{compat=1.15}
\usepackage{mathrsfs}


\title{\textbf{Filament 文献阅读 2\\Light propagation in field-ionizing media: Extreme nonlinear optics}\cite{geisslerLightPropagationFieldionizing1999}}                           %%%%% 标题
\author{骆禹天}                                           %%%%% 作者
\date{2024/**/**}                                        %%%%%% 日期

%%   定义一些数学格式
\newtheorem{theorem}{定理}[]
\newtheorem{definition}[theorem]{定义}
\newtheorem{lemma}[theorem]{引理}
\newtheorem{corollary}[theorem]{推论}
\newtheorem{example}[theorem]{例}
\newtheorem{proposition}[theorem]{命题}
\renewcommand{\abstractname}{\Large\textbf{摘要}}




%------------------------- 开始文档  -------------------------%



\begin{document}

\pagenumbering{Roman}         %%%%% 目录前页码统计按罗马数字

\maketitle

\setcounter{page}{1}          %%%%% 标题开始统计页数

%------------------------- 目录  ------------------------------%

\tableofcontents              %%%%% 插入目录

\newpage
\setcounter{page}{1}          %%%%% 目录后重新统计页数


\pagenumbering{arabic}        %%%%% 目录后统计页码为阿拉伯数字



%%%%%%%%%%%%%%%%%%%% 开始第一节  %%%%%%%%%%%%%%%%%%%%%



\section{Introduction}

对强场原子和等离子体现象的精确描述需要结合对介质中微观过程的分析来求解麦克斯韦方程。但这往往超出了当前计算机的能力，因而必须使用简化模型。在这封信中，作者推导了一种新颖的强场传播方程，该方程在以下两个方面简化和改进了现有模型:

（i）将非线性光学微扰范围中使用的概念扩展到强场范围。利用缓慢演化波 (SEW) 近似 [2]，导出了强场原子和等离子体物理的一阶传播 (FOP) 方程，该方程对于低密度介质和任意短脉冲持续时间有效。。

(ii) 先前宏观传播研究中使用的电离介质[8-10]的极化响应是不完整的[11]，并且不能正确描述基本过程，例如场电离引入的吸收损耗。作者的半经典分析揭示了一个额外的损失项。与氦气中的吸收测量进行比较，结果非常一致，并证实了在我们的实验中将吸收损失增加 3 倍的附加项的有效性和重要性。

\section{Derivation}

高强度脉冲在物质中的传播通过麦克斯韦方程结合介质 P 的偏振响应模型进行建模。只要坐标 (r x, y) 中的 P r, t 垂直于传播方向 (z ）在与中心波长 l0 相当的距离上表现出很小的变化，麦克斯韦方程可以用标量波动方程代替。在高斯单位中，对于线性极化电场，标量波动方程由下式给出
\begin{equation}
    \partial_z^2E\:+\:\nabla_\perp^2E\:-\:\frac1{c^2}\:\partial_t^2E\:=\:\frac{4\pi}{c^2}\:\partial_t^2P
    \label{eq:1}
\end{equation}
其中, c 是真空光速，≠z,t是对z和t的偏导数，=2是横向拉普拉斯算子。 FOP 方程的推导是在谱域中通过将 ansatz E r、v U r、v exp ivz c 插入到傅里叶变换方程中来进行的。 (1)，其中 E r, v R` 2` dt E r, t exp 2ivt 。






%--------------------  插入参考文献  --------------------%

\newpage

\clearpage           %%%  “目录”中显示的页码正确, 需要添加

\phantomsection      %%   “目录”中的链接能正确跳转, 需要添加 

\addcontentsline{toc}{section}{参考文献}   %向目录中添加条目, 以章的名义

\bibliography{ref}

\end{document}  